\section{Results}
\label{sec:results}

\subsection{What was implemented and what was tested}
\label{sec:coverage}

Eight parameter families were converted into MOSAICS decks and evaluated at
all-atom resolution: ff14SB \cite{maier2015ff14sb} and ff19SB
\cite{tian2020ff19sb} for protein; parmbsc0 \cite{perez2007bsc0}, parmbsc1
\cite{ivani2016bsc1}, OL15/OL3 \cite{zgarbova2015ol15,zgarbova2011ol3},
OL21/OL3 \cite{zgarbova2021ol21,zgarbova2011ol3} and OL24/OL3
\cite{zgarbova2025ol24,zgarbova2011ol3} for nucleic acid; and CHARMM36 for
nucleic acid \cite{hart2012charmm36dna,denning2011charmm36rna}. Across the
eleven systems they give the 26 system-and-force-field combinations of
Table~\ref{tab:matrix}, 54 MOSAICS-to-reference comparisons in all, 20 of the
26 cells carrying two or more references.

The cells were not all checked to the same strength, and the table says which
is which. The six original AMBER combinations on the heteroduplex and the
tetrapeptide were compared term by term against all three references, and the
ff14SB proteins and the OL21/OL3, OL15/OL3 and OL3 nucleic systems have a
GROMACS total as well as sander and OpenMM. Every ff19SB row is against sander
and OpenMM only, because the ParmEd route collapses its per-residue correction
maps onto one (F-06, Section~\ref{sec:methods-references}); the GROMACS cell is
blank, not failing. The CHARMM36 and OL24 rows are against OpenMM only.
CHARMM36 has one reference because its deck and its reference are both read
out of the same OpenMM build; on the Drew--Dickerson dodecamer the sander OL24
run read the full-precision restart coordinates where MOSAICS and OpenMM read
the three-decimal PDB, so it is a comparison of two geometries and is left out.
The 1efs CHARMM36 cell is the 816-atom $5'$-OH rebuild of the heteroduplex
(F-08) and is not comparable with the AMBER rows above it. The two 1F7Y rows
run on decks that exist only on an unmerged engine branch. Every run is a
single-point energy at one fixed geometry, in vacuum, with no cutoff and
dielectric~1.

\paragraph{Which chemistry the test systems reach}
A deck is a table, and a single-point energy exercises only the rows the
structure selects. The eleven systems reach all eight standard nucleic bases,
DA, DC, DG and DT on the dodecamer and A, C, G and U on the tetraloop and the
RNA 4-mer; $5'$-OH and $3'$-OH chain ends on five systems and $5'$-phosphate
interior chemistry on the padded heteroduplex; twenty amino-acid types across
ubiquitin, BPTI and the MHC heavy chain, including charged, aromatic,
histidine and proline residues, three disulfides, charged termini and two
capped fragments. What no system reaches is a $5'$-phosphate terminus under
any deck, an alternative protonation state, or more than one conformation: each
cell is one geometry, so a parameter that is wrong but inactive at that
geometry, or a torsion whose phase is wrong but degenerate there, would not
show. Finding F-21 is that failure mode caught in the act: an analytic force
that was identically zero passed every energy comparison on the peptide.

\subsection{Single-point agreement over 26 combinations}
\label{sec:matrix}

Table~\ref{tab:matrix} gives every total energy as each program returns it,
with the electrostatic term of MOSAICS and sander rescaled to the single
constant of Section~\ref{sec:constant}, and the difference of each reference
from MOSAICS. Figure~\ref{fig:energy-agreement} plots the same 54 differences,
and, where the engine output files resolve them, the differences term by term.

\begin{table}[tbp]
\centering
\footnotesize
\setlength{\tabcolsep}{4pt}
\caption{Single-point total energies on the eleven systems, in \kcal{}, vacuum,
no cutoff, every electrostatic term on the CODATA-derived constant
$332.0637133$. A dash means that reference was not run for that cell, for the
reason given in Section~\ref{sec:coverage}. The right-hand block is the
reference minus MOSAICS. The 1efs CHARMM36 row is the 816-atom $5'$-OH
rebuild of the heteroduplex, not the 818-atom structure of the rows above it;
the 1F7Y decks are on an unmerged engine branch. Source:
\texttt{work/preprint-I/energies/MATRIX.json}, and the engine output files it
names; the table is written from that file by \texttt{make\_tables.py}.}
\label{tab:matrix}
\resizebox{\textwidth}{!}{% generated by make_tables.py from energies/MATRIX.json
\begin{tabular}{@{}llrrrrrrr@{}}
\toprule
 & & \multicolumn{4}{c}{Total energy (kcal/mol)} & \multicolumn{3}{c}{Reference minus MOSAICS (kcal/mol)} \\
\cmidrule(lr){3-6}\cmidrule(lr){7-9}
System & Force field & MOSAICS & sander & OpenMM & GROMACS & sander & OpenMM & GROMACS \\
\midrule
1efs & OL21/OL3 & 2526.186378 & 2526.186713 & 2526.186575 & 2526.186580 & +0.000335 & +0.000197 & +0.000202 \\
 & OL15/OL3 & 2522.856302 & 2522.856613 & 2522.856501 & 2522.856506 & +0.000311 & +0.000199 & +0.000204 \\
 & OL24/OL3 & 2954.803590 & -- & 2954.803568 & -- & -- & -0.000022 & -- \\
 & CHARMM36 & 3775.995925 & -- & 3775.995927 & -- & -- & +0.000002 & -- \\
1BNA & OL21 & 2141.488688 & 2141.487481 & 2141.487431 & 2141.488017 & -0.001208 & -0.001258 & -0.000671 \\
 & OL15 & 2138.533139 & 2138.531881 & 2138.531881 & 2138.532468 & -0.001258 & -0.001258 & -0.000671 \\
 & CHARMM36 & 3509.022307 & -- & 3509.022309 & -- & -- & +0.000001 & -- \\
 & OL24 & 2954.122454 & -- & 2954.122861 & -- & -- & +0.000407 & -- \\
d(ACGT) & OL21 & 942.827346 & 942.827874 & 942.827764 & -- & +0.000528 & +0.000418 & -- \\
 & OL15 & 945.670885 & 945.671374 & 945.671303 & -- & +0.000489 & +0.000418 & -- \\
 & OL24 & 1080.088750 & -- & 1080.088749 & -- & -- & -0.000000 & -- \\
2KOC & OL3 & 347.450922 & 347.450964 & 347.451027 & 347.450987 & +0.000042 & +0.000106 & +0.000065 \\
 & CHARMM36 & 2106.624466 & -- & 2106.624466 & -- & -- & +0.000000 & -- \\
r(ACGU) & OL3 & 2584.335888 & 2584.336025 & 2584.335946 & -- & +0.000137 & +0.000058 & -- \\
1UBQ & ff14SB & -546.960403 & -546.959674 & -546.959594 & -546.959625 & +0.000729 & +0.000809 & +0.000778 \\
 & ff19SB & -1022.050252 & -1022.050874 & -1022.050802 & -- & -0.000622 & -0.000550 & -- \\
5PTI & ff14SB & -332.312157 & -332.311788 & -332.311563 & -332.311590 & +0.000369 & +0.000594 & +0.000567 \\
 & ff19SB & -688.753926 & -688.711488 & -688.711342 & -- & +0.042438 & +0.042584 & -- \\
ACE-AGSV-NME & ff14SB & 96.262476 & 96.262430 & 96.262484 & 96.262481 & -0.000046 & +0.000008 & +0.000005 \\
 & ff19SB & 69.244797 & 69.244830 & 69.244805 & -- & +0.000033 & +0.000008 & -- \\
Tax 11--19 & ff14SB & -23.358508 & -23.358422 & -23.358416 & -- & +0.000087 & +0.000092 & -- \\
 & ff19SB & -82.278687 & -82.278622 & -82.278606 & -- & +0.000066 & +0.000081 & -- \\
1HHK & ff14SB & -475.878441 & -475.874656 & -475.873684 & -- & +0.003785 & +0.004757 & -- \\
 & ff19SB & -2937.222999 & -2937.219256 & -2937.218706 & -- & +0.003743 & +0.004293 & -- \\
1F7Y & ff14SB/OL3 & 14238.328640 & 14238.329782 & 14238.329881 & -- & +0.001142 & +0.001241 & -- \\
 & ff19SB/OL3 & 13685.351529 & 13685.353082 & 13685.353068 & -- & +0.001553 & +0.001539 & -- \\
\bottomrule
\end{tabular}
}
\end{table}

\begin{figure}[tbp]
\centering
\includegraphics[width=\linewidth]{energy-agreement.pdf}
\caption{Energy agreement over the 26 combinations. (a) One marker per
MOSAICS-to-reference comparison of the total energy, 54 in all, grouped by
system class; the vertical axis is $|E_{\mathrm{MOSAICS}} -
E_{\mathrm{reference}}|$ on a log scale. Marker shape and colour identify the
reference engine. (b) The same comparisons resolved into the terms the engines
share, for the cells whose output files carry the split; where an OpenMM file
prints only a combined non-bonded energy the comparison appears in the last
column. Hollow markers fall below the print resolution of the reference file
they were read from and are print quantisation, not measured disagreement.
Dashed line, $kT$ at 300~K, 0.596~\kcal{}; dotted line, 0.001~\kcal{}. Every
value is re-derived from \texttt{MATRIX.json} and the engine outputs it names
by \texttt{figures/energy-agreement.py}.}
\label{fig:energy-agreement}
\end{figure}

Nineteen of the 26 cells agree with every reference available to them to
within 0.001~\kcal{}, and thirteen to within 0.0004. Over the 26 cells the
median of the largest gap to any reference is 0.00037~\kcal{}; the smallest is
$3.2\times10^{-7}$, on the tetraloop under CHARMM36, and the largest is
0.0426, on BPTI under ff19SB. For scale, $kT$ at 300~K is 0.596~\kcal{}, and the
sander column can resolve nothing below about $2\times10^{-4}$. The seven cells
outside 0.001~\kcal{} are, in order, BPTI ff19SB (0.0426), the two
peptide--MHC cells (0.0048 and 0.0043), the two S15--rRNA cells (0.0012 and
0.0016) and the dodecamer under OL21 and OL15 (0.0013 each). Each has a
mechanism, and they are of three kinds.

\paragraph{An improper read in two atom orders}
The BPTI ff19SB cell is 0.0426~\kcal{} below the references on the dihedral
term and nowhere else: bond, angle, Lennard-Jones and electrostatics carry the
same residuals as the ff14SB cell on the same coordinates, and the correction
map matches OpenMM to $5.8\times10^{-13}$. Splitting the dihedral term puts
the whole gap in the impropers, 5.5085 against 5.5506~\kcal{}. It is one
improper, at the backbone nitrogen of ASP~3. The AMBER topology writes that
improper in the order C--CA--N--H and the MOSAICS residue template in the order
$-$C--H--N--CA; the two orders define different dihedral angles, which coincide
when the nitrogen is planar and part company when it is not, and at ASP~3 in
the deposited coordinates they read $-156.8^{\circ}$ and $+158.4^{\circ}$,
some $22^{\circ}$ from planar. Evaluating that one improper in the topology's
order, with the two conventions for the phase and the deck's coefficient
rounding accounted for, reproduces the whole gap to $8.5\times10^{-8}$~\kcal{},
and the accounting is committed beside the run. The same mechanism produced a
0.46~\kcal{} gap at LEU~2 of the Tax nonamer on an earlier build of that
structure, which is why the peptide's first amide hydrogen is now placed in the
peptide plane by the engine's builder and why the nonamer cells in
Table~\ref{tab:matrix} agree to $9\times10^{-5}$. The ff14SB cell on the same
BPTI coordinates is within $6\times10^{-4}$, so the parameters are not in
question: the improper order is.

\paragraph{Six adenine impropers on the dodecamer}
Under OL21 and OL15 the dodecamer sits 0.0013~\kcal{} above sander and OpenMM
and 0.0007 above GROMACS. Resolved into terms the excess is on the torsion
term, $+1.5\times10^{-3}$, of which $+1.0\times10^{-3}$ is improper, and it
was located rather than guessed. \texttt{tleap} under DNA.OL21 writes the
adenine N9 improper as C8/C4/N9/C1$'$ and under ff99bsc0 as C4/C8/N9/C1$'$;
swapping the first two atoms is not a symmetry of the term when the group is
not exactly planar; MOSAICS agrees with the parmbsc0 order to eight figures,
and the six instances that differ, on DA~5, 16 and 17, account for the whole
$1.0\times10^{-3}$. The same term on the heteroduplex, whose entire improper
energy is $5.3\times10^{-4}$~\kcal{}, is within $4.3\times10^{-5}$, so the
difference between the two structures is the geometry the impropers are
evaluated at, not the parameters.

\paragraph{Two large systems not yet resolved into terms}
The peptide--MHC complex, 6,139 atoms, is 0.0038~\kcal{} above sander and 0.0048
above OpenMM under ff14SB, and 0.0037 and 0.0043 under ff19SB, out of a total of
$-476$ and $-2937$; per atom that is below $10^{-6}$~\kcal{}, and the two
force fields give the same offset within $5\times10^{-4}$, which points at the
non-bonded or angle terms they share rather than at the backbone terms they do
not. The S15--rRNA complex, 3,321 atoms, is 0.0012 and 0.0016~\kcal{} below the
references, and its terms are resolved: bond agrees to $2\times10^{-11}$, the
correction map to $2\times10^{-11}$, torsion and improper to $4\times10^{-4}$
and $8\times10^{-5}$, while the angle term carries $-1.55\times10^{-3}$ and the
non-bonded sum $-3.7\times10^{-3}$ in the two directions that nearly cancel.
Those two systems are the only ones whose reference coordinates were written
from the three-decimal PDB rather than from \texttt{tleap}'s restart, and the
S15--rRNA decks are on a branch; both cells are reported as they stand and
neither is claimed to be closed.

\paragraph{The angle term, on every system}
The one residual that appears on every system whose terms were resolved is on
the angle term, and it always has the same sign. MOSAICS minus OpenMM on that
term is $-1.76\times10^{-4}$~\kcal{} on the heteroduplex, $-9.6\times10^{-6}$
on the tetrapeptide, $-1.0\times10^{-4}$ on the tetraloop, $-3.4\times10^{-4}$
on the dodecamer, $-5.1\times10^{-4}$ on BPTI and $-6.5\times10^{-4}$ on
ubiquitin, growing with the number of angles. It is not a MOSAICS error. The
AMBER topology builder writes every equilibrium angle as the published value
multiplied by a seven-digit approximation of $\pi/180$, while MOSAICS reads the
value in degrees and converts it exactly, so the reference carries the rounding
(F-12). It is reported and not corrected, and it is the largest residual on
every AMBER cell that is inside 0.001~\kcal{}.

\subsection{Term by term on the two reference systems}
\label{sec:energies}

Before the numbers, the resolution of the files they are read from, because
several of the differences below are smaller than the output they come out of.
sander prints four decimals, an absolute floor near $2\times10^{-4}$~\kcal{}.
GROMACS energies are read from \texttt{sp.xvg}, which carries six decimals in
kJ/mol, a floor of $2.4\times10^{-7}$~\kcal{}. OpenMM's single-point output
carries twelve decimal places, a floor of about $5\times10^{-13}$~\kcal{}.
MOSAICS prints sixteen significant digits. A difference smaller than the floor
of the file it is taken from is the last printed digit of that file and not a
measurement, and no such number is quoted here as a residual.

On every one of the six original AMBER combinations, the largest difference
between MOSAICS and OpenMM on any energy term is $1.8\times10^{-4}$~\kcal{} for
the four heteroduplex decks and $9.6\times10^{-6}$~\kcal{} for the two peptide
decks, in both cases on the angle term and for the reason just given. Every
other term on every other row is smaller than this. The margin varies widely:
the parmbsc1 dihedral residual, $+1.03\times10^{-4}$~\kcal{}, is within a factor
of two of it, while the two peptide bond terms agree past the twelfth decimal
place OpenMM writes, which is as far as the comparison can be taken and no
further. Table~\ref{tab:ol21} gives the full four-engine breakdown for one force
field, OL21/OL3 on 1efs, so that the raw numbers are on the page, and
Table~\ref{tab:summary} gives MOSAICS minus each reference, per term, for all
six combinations.

\begin{table}[t]
\centering
\small
\setlength{\tabcolsep}{5pt}
\caption{OL21/OL3 on 1efs (818 atoms, net charge $-26$, vacuum, no cutoff), in
\kcal{}, as each engine prints them. The dihedral row includes impropers, and the
Lennard-Jones and electrostatic rows include the 1--4 terms. sander gives four
decimals (floor $2\times10^{-4}$~\kcal{}), the GROMACS column is converted from
\texttt{sp.xvg} values carrying six decimals in kJ/mol (floor
$2.4\times10^{-7}$), OpenMM writes twelve decimals (floor $5\times10^{-13}$)
and MOSAICS writes sixteen significant digits. \emph{corrected} means the
electrostatic term rescaled to the single constant of
Section~\ref{sec:constant}; the MOSAICS column is divided by the unrounded
product $0.529177\times627.50921$ and the sander column by
$18.2223^{2}=332.05221729$, so either corrected entry can be reproduced digit
for digit from the printed one above it. Sources:
\texttt{evidence/1efs/ol21\_ol3/\{mosaics,sander,openmm,gromacs\}/}.}
\label{tab:ol21}
\begin{tabular}{@{}lrrrr@{}}
\toprule
Term & MOSAICS & sander & OpenMM & GROMACS \\
\midrule
bond                   &   70.99877741 &   70.9988    &   70.99877741 &   70.99877749 \\
angle                  &  220.56489609 &  220.5651    &  220.56507252 &  220.56507290 \\
dihedral               &  596.17955421 &  596.1796    &  596.17959679 &  596.17959369 \\
Lennard-Jones          & $-$219.29643681 & $-$219.2964 & $-$219.29643710 & $-$219.29646343 \\
electrostatic, printed & 1857.73806648 & 1857.6753    & 1857.73956494 & 1857.73959895 \\
electrostatic, corrected & 1857.73958864 & 1857.73961475 & 1857.73956494 & 1857.73959895 \\
\midrule
total, corrected       & 2526.18637954 & 2526.18671475 & 2526.18657455 & 2526.18657959 \\
\bottomrule
\end{tabular}
\end{table}

\begin{table}[t]
\centering
\footnotesize
\setlength{\tabcolsep}{2.5pt}
\caption{MOSAICS minus each reference engine, per term, in \kcal{}, on the six
original combinations. Electrostatics are corrected on both sides to
$332.0637133$~kcal\,\AA\,mol$^{-1}$\,$e^{-2}$ (Section~\ref{sec:constant})
before the difference is taken, with the MOSAICS side divided by the unrounded
$0.529177\times627.50921 = 332.06344122017$. \emph{n.r.} means the difference
is smaller than the resolution of that reference's own output file, so the
value in the committed evidence is a print artefact; the 27 censored sander
entries run from $2.0\times10^{-6}$ to $5.1\times10^{-5}$ in the raw files.
The GROMACS CMAP entry is the conversion-route limit F-06 and not a
disagreement about physics. \emph{--} marks a term the force field does not
have. Source: \texttt{evidence/results.txt} and the per-combination output
files under \texttt{evidence/}.}
\label{tab:summary}
\begin{tabular}{@{}lrrrrrr@{}}
\toprule
Term & OL21/OL3 & OL15/OL3 & parmbsc1 & parmbsc0 & ff14SB & ff19SB \\
\midrule
\multicolumn{7}{@{}l}{\emph{Reference: OpenMM 8.5.2, output floor $5\times10^{-13}$~\kcal{}}}\\
bond      & n.r. & $+3.78\times10^{-9}$ & $-8.30\times10^{-8}$ & $-8.30\times10^{-8}$ & n.r. & n.r. \\
angle     & $-1.76\times10^{-4}$ & $-1.76\times10^{-4}$ & $-1.76\times10^{-4}$ & $-1.76\times10^{-4}$ & $-9.57\times10^{-6}$ & $-9.57\times10^{-6}$ \\
dihedral  & $-4.26\times10^{-5}$ & $-4.45\times10^{-5}$ & $+1.03\times10^{-4}$ & $+3.59\times10^{-5}$ & $+1.27\times10^{-6}$ & $+1.28\times10^{-6}$ \\
cmap      & --  & --  & --  & --  & --  & $-5.69\times10^{-8}$ \\
Lennard-Jones & $+2.89\times10^{-7}$ & $+2.71\times10^{-7}$ & $+3.81\times10^{-6}$ & $+3.81\times10^{-6}$ & $+1.15\times10^{-8}$ & $+1.15\times10^{-8}$ \\
electrostatic & $+2.37\times10^{-5}$ & $+2.37\times10^{-5}$ & $+2.37\times10^{-5}$ & $+2.37\times10^{-5}$ & $+4.18\times10^{-7}$ & $+4.18\times10^{-7}$ \\
\addlinespace
\multicolumn{7}{@{}l}{\emph{Reference: sander 24.0, output floor $2\times10^{-4}$~\kcal{}}}\\
bond      & n.r. & n.r. & n.r. & n.r. & n.r. & n.r. \\
angle     & $-2.04\times10^{-4}$ & $-2.04\times10^{-4}$ & $-2.04\times10^{-4}$ & $-2.04\times10^{-4}$ & n.r. & n.r. \\
dihedral  & n.r. & n.r. & n.r. & n.r. & n.r. & n.r. \\
cmap      & --  & --  & --  & --  & --  & n.r. \\
Lennard-Jones & n.r. & n.r. & n.r. & n.r. & n.r. & n.r. \\
electrostatic & n.r. & n.r. & n.r. & n.r. & n.r. & n.r. \\
\addlinespace
\multicolumn{7}{@{}l}{\emph{Reference: GROMACS 2025.4, output floor $2.4\times10^{-7}$~\kcal{}}}\\
bond      & n.r. & n.r. & n.r. & n.r. & n.r. & n.r. \\
angle     & $-1.77\times10^{-4}$ & $-1.77\times10^{-4}$ & $-1.77\times10^{-4}$ & $-1.77\times10^{-4}$ & $-9.60\times10^{-6}$ & $-9.60\times10^{-6}$ \\
dihedral  & $-3.95\times10^{-5}$ & $-4.16\times10^{-5}$ & $+1.07\times10^{-4}$ & $+3.90\times10^{-5}$ & $+1.30\times10^{-6}$ & $+1.38\times10^{-6}$ \\
cmap      & --  & --  & --  & --  & --  & $-3.02\times10^{-1}$ (F-06) \\
Lennard-Jones & $+2.66\times10^{-5}$ & $+2.66\times10^{-5}$ & $+3.01\times10^{-5}$ & $+3.01\times10^{-5}$ & $+3.02\times10^{-6}$ & $+3.02\times10^{-6}$ \\
electrostatic & $-1.03\times10^{-5}$ & $-1.03\times10^{-5}$ & $-1.03\times10^{-5}$ & $-1.03\times10^{-5}$ & $+6.49\times10^{-7}$ & $+6.49\times10^{-7}$ \\
\bottomrule
\end{tabular}
\end{table}

\paragraph{On the two non-bonded terms, the reference engines disagree with each
other more than MOSAICS disagrees with either}
OpenMM and GROMACS are independently written codes, and on 1efs they were given
the same parameters and the same coordinates. They differ from one another by
$+2.63\times10^{-5}$~\kcal{} on Lennard-Jones and $-3.40\times10^{-5}$ on
electrostatics, on all four nucleic decks. MOSAICS minus OpenMM on the same two
terms is $+2.89\times10^{-7}$ on Lennard-Jones under OL21/OL3
($+3.81\times10^{-6}$ under the two bsc decks) and $+2.20\times10^{-5}$ on
electrostatics, the latter with both sides corrected to OpenMM's own printed
constant. That pattern, bonded terms with few parameters surviving and
non-bonded terms with thousands not, is what a text-format topology conversion
would produce, since a GROMACS \texttt{.top} stores fewer digits per parameter
than the binary topology the other two engines read. We have not tested that
attribution directly. What the numbers establish is a bound that has nothing to
do with MOSAICS: two mature MD codes, handed the same force field through the
ordinary conversion route, land $3\times10^{-5}$~\kcal{} apart on an 818-atom
system, and on the two non-bonded terms MOSAICS sits below that bound against
either reference.

\paragraph{The peptide test is a targeted test of what ff19SB adds}
Between ff14SB and ff19SB on the same capped peptide, the bond, angle,
Lennard-Jones and electrostatic terms are bit-identical in all four engines.
MOSAICS returns the same bond (0.14749181), bend (0.98732525), Lennard-Jones
(134.50567845) and electrostatic ($-71.35989834$~\kcal{}) energies in both
runs, agreeing to the last of the sixteen digits it prints. The whole ff14SB-to-ff19SB difference is
$-27.01767928$~\kcal{}, of which $-25.09643922$ is the backbone torsion and
$-1.92124006$ is the correction map; every other term contributes exactly zero.
The ff19SB row is therefore the correction-map machinery applied on top of an
already-verified ff14SB, tested with everything else held fixed, and four
residues select four of the sixteen maps the deck loads. On ubiquitin, BPTI and
the two complexes the same machinery is exercised on every amino-acid type the
deck holds, and the correction-map term matches OpenMM to between
$3\times10^{-13}$ and $2\times10^{-11}$~\kcal{} on all six ff19SB cells.

\paragraph{The nucleic decks are different parameter sets, and each is
reproduced as it stands}
All three references return the same bond, angle, Lennard-Jones and
electrostatic energies for all four AMBER nucleic decks on 1efs; only the
dihedral term moves. That is the design of those decks showing through, and it
makes the dihedral row of Table~\ref{tab:summary} the row that carries the
result: MOSAICS follows each of the four dihedral values to between
$3.6\times10^{-5}$ and $1.0\times10^{-4}$~\kcal{} against OpenMM and GROMACS,
without being retuned for any of them. On the tetraloop the OL15/OL3 and OL21/OL3
terminal decks return the same energy, 347.4514261182977~\kcal{} before the
constant correction, to all sixteen printed digits on every term, because OL15
and OL21 touch DNA torsions only and a pure RNA chain never reaches them;
Table~\ref{tab:matrix} carries that one measurement once, as OL3. The
size of the gap between two decks' absolute torsion energies is not reported
here and is not a quantity this protocol can produce.

\paragraph{CHARMM36 nucleic acid, one reference}
An OpenMM reference for the CHARMM36 nucleic deck was built from
\texttt{charmm36\_2024.xml} and evaluated on the two free-ended 4-mers, the
dodecamer, the tetraloop and the $5'$-OH rebuild of the heteroduplex. The deck
covers 24 residues (every base as a $5'$, internal and $3'$ unit, DNA and RNA)
with 237 bonds, 562 bends, 1036 torsion entries and 15\,753 non-bonded pairs.
Table~\ref{tab:charmm36} gives the two 4-mers term by term. Every bonded term
matches to the digits OpenMM prints, and that includes the two CHARMM terms the
AMBER force fields do not have, the Urey--Bradley 1--3 harmonics and the
harmonic impropers. What the improper row shows needs care: its total is
$1.0\times10^{-3}$~\kcal{} over 188 terms, so the structure sits at the improper
minimum for practical purposes, and that is exactly where the engine's truncated
sixth-order expansion and CHARMM's quadratic coincide by construction.
Ten-digit agreement there confirms that the 188 entries were converted and
selected correctly; it does not exercise the functional form away from
equilibrium. The electrostatic term is the only one that does not match as
printed, at $9.8\times10^{-5}$~\kcal{} on the RNA 4-mer and $2.9\times10^{-4}$
on the DNA 4-mer, and the charge-to-energy constant accounts for that offset to
four significant figures: corrected exactly as the AMBER rows are corrected, the
two residuals fall to $-2.4\times10^{-10}$ and $-9.0\times10^{-10}$~\kcal{}. On
the three larger structures the corrected totals differ from OpenMM by
$-3.5\times10^{-7}$ (dodecamer), $+3.2\times10^{-7}$ (tetraloop) and
$-2.1\times10^{-6}$~\kcal{} (816-atom heteroduplex). That is the closest
agreement anywhere in this work, and Section~\ref{sec:constantgeneral} says
why.

\begin{table}[t]
\centering
\small
\caption{CHARMM36 nucleic acid against OpenMM 8.5.2, in \kcal{}, on the two
free-ended 4-mers with 5TER/3TER termini in vacuum, as printed. r(ACGU) is 127
atoms and d(ACGT) is 126. Urey--Bradley is CHARMM's 1--3 harmonic term, which in
OpenMM sits inside \texttt{HarmonicBondForce} and is separated here. The two
d(ACGT) bonded columns are inflated by a defect in that committed structure, in
which one H--C--H angle is $3.15^{\circ}$ instead of $109.5^{\circ}$ (F-17,
Section~\ref{sec:limits}); the comparison between the two programs is
unaffected, because both read the same coordinates. The electrostatic rows
differ by the charge-to-energy constant and by little else. Sources:
\texttt{evidence/charmm36\_rna\_acgu.tsv}, \texttt{evidence/charmm36\_dna\_acgt.tsv}.}
\label{tab:charmm36}
\begin{tabular}{@{}lrrrr@{}}
\toprule
& \multicolumn{2}{c}{r(ACGU)} & \multicolumn{2}{c}{d(ACGT)} \\
\cmidrule(lr){2-3}\cmidrule(lr){4-5}
Term & MOSAICS & OpenMM & MOSAICS & OpenMM \\
\midrule
bond           &   28.51501806 &   28.51501806 &   18.92184909 &   18.92184909 \\
Urey--Bradley  &   27.33939420 &   27.33939420 &  167.65137978 &  167.65137978 \\
angle          &  166.36120948 &  166.36120948 &  412.78809431 &  412.78809431 \\
torsion        &  114.02952401 &  114.02952401 &  122.94661546 &  122.94661546 \\
improper       &    0.00102252 &    0.00102252 &    0.00102555 &    0.00102555 \\
1--4 Lennard-Jones & 230.89575106 & 230.89575106 & 1412.27140098 & 1412.27140098 \\
Lennard-Jones  & 2564.74970076 & 2564.74970076 &  176.68820406 &  176.68820406 \\
electrostatic  & $-$119.98431509 & $-$119.98441340 & $-$356.57731681 & $-$356.57760898 \\
\midrule
total          & 3011.90730500 & 3011.90720669 & 1954.69125242 & 1954.69096025 \\
\bottomrule
\end{tabular}
\end{table}

\subsection{The charge-to-energy constant, and why every electrostatic number is corrected}
\label{sec:constant}

The four programs do not use the same constant to turn charges into an energy.
This was measured and not assumed, by the two-charge topology of
Section~\ref{sec:methods-coulomb}, and Table~\ref{tab:coulomb} gives the four
values. Every electrostatic term in this work is put on one target constant
before any difference is taken, and that target has to be a number with a
derivation, or the correction becomes a preference. It is
\begin{equation*}
  \frac{e^2 N_{\mathrm A}}{4\pi\varepsilon_0 \times 1\,\text{\AA} \times
  4.184\,\mathrm{J\,cal^{-1}}}
  = 332.0637133\ \mathrm{kcal}\,\text{\AA}\,\mathrm{mol}^{-1}\,e^{-2},
\end{equation*}
evaluated with the CODATA 2018 recommended values\cite{codata2018}. The
elementary charge, $e = 1.602176634\times10^{-19}$~C, and the Avogadro
constant, $N_{\mathrm A} = 6.02214076\times10^{23}$~mol$^{-1}$, are both exact
by the 2019 SI definitions; $\varepsilon_0 = 8.8541878128\times10^{-12}$
F\,m$^{-1}$; and the thermochemical calorie is exactly 4.184~J. The 2022 revision of $\varepsilon_0$ gives $332.0637131$
instead, lower by $6.8\times10^{-10}$ relative, which is worth
$1.3\times10^{-6}$~\kcal{} on the 1efs electrostatic term.

\begin{table}[t]
\centering
\small
\caption{The charge-to-energy constant of each program, in
kcal\,\AA\,mol$^{-1}$\,e$^{-2}$. sander, OpenMM and GROMACS were measured with
two unit charges 1~\AA{} apart in vacuum, everything else switched off; the
MOSAICS entry is evaluated from the constants compiled into the engine,
$0.529177 \times 627.50921$. sander's measured value is limited by its
four-decimal print precision; its exact value is $18.2223^2 = 332.05221729$,
and that full value is what corrects the sander column. The OpenMM and GROMACS
measurements were recorded to six decimals, which bounds their relative
difference from the target at $1.5\times10^{-9}$ and says nothing finer.
Source: \texttt{evidence/coulomb-constants/}.}
\label{tab:coulomb}
\begin{tabular}{@{}lrrr@{}}
\toprule
Engine & Constant & Relative to sander & Relative to $332.0637133$ \\
\midrule
sander  & 332.052200 & $0$                   & $-3.5\times10^{-5}$ \\
OpenMM  & 332.063713 & $+3.467316\times10^{-5}$ & $<1.5\times10^{-9}$, print-limited \\
GROMACS & 332.063713 & $+3.467284\times10^{-5}$ & $<1.5\times10^{-9}$, print-limited \\
MOSAICS & 332.063441 & $+3.385377\times10^{-5}$ & $-8.2\times10^{-7}$ \\
\bottomrule
\end{tabular}
\end{table}

sander is the outlier, not MOSAICS. It stores topology charges pre-multiplied
by 18.2223 and its constant is that number squared, which is an older value;
OpenMM and GROMACS, written independently of each other, agree with each other
and with the derived target to every decimal their measurement was recorded to.
MOSAICS sits $8.2\times10^{-7}$ below the target because three constants
compiled into the engine are rounded (F-13). The size of this effect exceeds
every difference the comparison is trying to resolve. On 1efs the sander
electrostatic term moves by $+0.0643147$~\kcal{} when rescaled, and the MOSAICS
term by $+1.522158\times10^{-3}$; the sander shift is roughly 360 times the
largest residual left after correction and about a tenth of $kT$ at 300~K. An
uncorrected cross-engine comparison of an 818-atom charged system would
therefore report a disagreement between sander and OpenMM bigger than
everything else on the page, and it would be reporting a unit convention.

Which constant is the target has to be stated, because the choice is worth
something. Corrected to OpenMM's own printed constant instead of the derived
one, MOSAICS minus OpenMM on electrostatics is $+2.202\times10^{-5}$~\kcal{} on
all four 1efs rows against the $+2.370\times10^{-5}$ of
Table~\ref{tab:summary}; the $1.68\times10^{-6}$ between the two frames is what
a six-decimal printed constant costs, and neither choice changes any
conclusion. After correction, $+2.37\times10^{-5}$~\kcal{} remains on the 1efs
electrostatic term, $1.3\times10^{-8}$ of the term, and we cannot account for
it. One hypothesis fits what is on the page. The AMBER decks reach MOSAICS
through a text route, and a prmtop writes charges in \texttt{\%FORMAT(5E16.8)},
about $10^{-8}$ relative per charge, the same order as the residual; the
CHARMM36 deck does not take that route, its charges are read out of a built
OpenMM \texttt{System} at full double precision, and its corrected residuals
are three orders of magnitude smaller on the same engine with the same
correction. Testing this needs a diff of the deck charges against the prmtop
charges at full precision, which is one command and is not in the evidence.

\subsection{Forces}
\label{sec:forces}

Analytic forces were checked against a central difference of the same energy,
term by term, over a ten-point step sweep from $10^{-2}$ to $3\times10^{-7}$
Bohr, on six Cartesian systems and, against a staging coordinate, on four
staged systems: 61 sweeps in all, covering five terms and their sum on every
system, plus the correction map on the ff19SB peptide. The residual is the
largest $|$numerical $-$ analytic$|$ over every atom and every component at a
given step. Figure~\ref{fig:force-convergence} plots the total-energy sweep on
six of the systems.

\begin{figure}[tbp]
\centering
\includegraphics[width=\linewidth]{force-convergence.pdf}
\caption{Convergence of the central-difference check on the total-energy
gradient for six of the 61 sweeps, one per Cartesian system and force field.
Each point is the largest $|$numerical $-$ analytic$|$ force component over all
atoms at that step. The grey line has slope exactly 2 through the coarsest
point; the printed exponent is the least-squares fit over the fixed window
$10^{-2}$ to $3\times10^{-4}$~Bohr. Below the window each sweep leaves the
falling branch where round-off in the energy overtakes the discretisation
error, which is a floor and not a disagreement. Source:
\texttt{evidence/forces/*.all.tsv}.}
\label{fig:force-convergence}
\end{figure}

A residual at a given step is not an error in the force. It is dominated by the
discretisation error of the difference at coarse steps and by round-off in the
energy evaluation at fine ones, and it is smallest where the two cross. At the
coarsest step, $10^{-2}$~Bohr, residuals across the set run up to
$3.6\times10^{-3}$~Ha/Bohr, which is what a central difference of that width
is expected to give. The number worth quoting is the minimum over each sweep.
The largest of the 61 minima is $3.8\times10^{-9}$~Ha/Bohr, on the non-bonded
staging gradient of the RNA 4-mer; the smallest is $4.5\times10^{-13}$, on the
peptide torsion under ff19SB. For scale, 0.002~\kcal{} per Bohr is
$3.2\times10^{-6}$~Ha/Bohr, and every one of the 61 minima is smaller than that
by a factor of at least 800.

The exponent is the measurement, and it was fitted over a step window fixed in
advance instead of chosen per sweep: $10^{-2}$ to $3\times10^{-4}$~Bohr, the
four coarsest steps, which every one of the 61 sweeps spans. Over that window
the fitted exponents run from 1.994 to 2.009, with a mean of 2.000 and a
standard deviation of 0.002. The fits are tight: the largest root-mean-square
residual of any of the 61 is 0.003 in $\log_{10}$ of the force residual, and no
single point departs from its own fit by more than 0.005. That is the
second-order convergence a central difference gives when the analytic force
really is the derivative of the energy. Below the window every sweep leaves the
falling branch at its own step and flattens, which is round-off and not
disagreement; the correction map leaves it earliest and climbs again, reaching
$1.6\times10^{-5}$~Ha/Bohr at $3\times10^{-7}$~Bohr, because the map is
interpolated on a grid and is not smooth to machine precision.

The sweep found one defect that no energy comparison in this paper could have
found. The correction-map analytic force was identically zero at a dihedral of
exactly $\pi$, which is where a peptide built extended sits; it is fixed and
pinned by a test (F-21). The energy was never affected, and every sweep
reported here is post-repair.

\subsection{What one Monte Carlo step costs}
\label{sec:cost}

One Monte Carlo step on the 818-atom heteroduplex under OL21/OL3 costs
3.46~ms on one core of an Apple M5 laptop, with 449~ms of setup that happens
once per run; the linear fit over $N = 0$ to $100$ steps has $R^2 = 0.99968$
(Figure~\ref{fig:cost-fit}). A step is one torsional proposal plus one full
energy evaluation over every non-excluded pair, so the figure is an upper bound
on the energy evaluation alone; the proposal is inside it and is not separated
here. It is a cost, not a benchmark: nothing here compares MOSAICS with another
program on the same hardware, the engine has never been tuned for the speed of
a single evaluation, and the number says nothing about scaling or about other
sizes. What it settles is the question a reader of the 26-cell matrix is
entitled to ask next, which is what one of those evaluations costs when the
engine is actually sampling. The strength the method claims lies elsewhere, in
the number of evaluations a guided proposal needs to cross a barrier that
molecular dynamics must climb in femtosecond steps, and that ratio is not
measured in this paper.

\begin{figure}[tbp]
\centering
\includegraphics[width=0.5\linewidth]{cost-fit.pdf}
\caption{Wall clock against number of Monte Carlo steps for the committed
1efs OL21/OL3 input, 818 atoms, serial, one core of an Apple M5. Small hollow
markers are the five repeats at each $N$, filled markers their medians, the
line the least-squares fit to the medians. The one 866~ms outlier at $N=0$ is
why medians are used. Source:
\texttt{work/preprint-I/runs/2026-08-31-cost-818/result.json}, re-fitted by
\texttt{figures/cost-fit.py}, which exits if its fit disagrees with the
committed one.}
\label{fig:cost-fit}
\end{figure}

\subsection{\studio{}: the same engine, verified in the browser}
\label{sec:studio-results}

The Studio is only as good as the engine it carries, so its verification is of
the same kind as the rest of this paper: the same input, run under the native
engine and under the WebAssembly one, compared number by number.
Figure~\ref{fig:studio} shows the page at two of its four steps on the run used
for that check.

\begin{figure}[p]
\centering
\begin{subfigure}[t]{0.84\linewidth}
\includegraphics[width=\linewidth]{studio-prepare.png}
\caption{}
\label{fig:studio-prepare}
\end{subfigure}
\par\smallskip
\begin{subfigure}[t]{0.84\linewidth}
\includegraphics[width=\linewidth]{studio-understand.png}
\caption{}
\label{fig:studio-understand}
\end{subfigure}
\caption{\studio{} in the browser, version 1.2.6, on the CGCGCG duplex used
for the cross-check of Table~\ref{tab:studio}. (a) The Prepare step: the built
structure, 376 atoms in 12 residues across two chains with real $5'$ and $3'$
ends; the force field chosen with its reason, here the OL21/OL3 terminal deck
because the chains carry real ends; the validation gate reporting that the
ordinary bonded and non-bonded energies for every residue type present
reproduce a committed reference; and the 3Dmol.js viewer. The left rail names
the engine as WebAssembly running inside the tab, with the Studio version,
engine commit, engine digest and build date. (b) The Understand step after the
BFGS minimisation: initial and final energies, the number of energies
reported, the wall time in this tab and the exit code; energy along the run;
RMSD to the first frame along the trajectory; and the final structure. The
energy-against-RMSD plot, the residue-by-residue comparison, the graph export
and the record follow further down the page. The Start and Run steps are in
Supporting Information Figure~S1.}
\label{fig:studio}
\end{figure}

\paragraph{Native against WebAssembly}
The Studio's documented case study builds CGCGCG as a B-DNA duplex, minimises
it with the \texttt{local-minimum} preset (BFGS to a gradient tolerance of
$10^{-7}$) on the OL21/OL3 terminal deck, exports a residue graph and writes
the report. Replaying the deck the Studio writes under the native Apple-silicon
engine and under the WebAssembly engine gives Table~\ref{tab:studio}: the same
exit code, the same 2,040 BFGS iterations, the same first and last energies to
the two decimals the engine's energy file carries, a largest difference along
the series of 0.25~\kcal{}, and a final-structure RMSD of 0.0045~\AA{} between
the two. A headless test drives the whole pipeline, the same Python archive and
the same bridge the page uses, and fails unless those numbers reproduce. The
run in Figure~\ref{fig:studio} is that case study executed on the live page in
a headless Chromium session on 8 September 2026: 2,040 energies reported,
1347.94 to $-1102.39$~\kcal{}, 25.1~s of wall time in the tab against 20.8~s
under Node and 11.5~s for the native engine.

\begin{table}[t]
\centering
\small
\caption{The CGCGCG case study replayed under the native engine and under the
WebAssembly engine, from the Studio's build record. The native binary is the
Apple-silicon 3.9.1 build at engine commit \texttt{ea00b02}; the WebAssembly
build is from commit \texttt{73f598a} (7 September 2026). Source:
\texttt{web/engine/BUILD.md} in the Studio repository.}
\label{tab:studio}
\begin{tabular}{@{}lrr@{}}
\toprule
 & native, arm64 & WebAssembly \\
\midrule
exit code & 0 & 0 \\
BFGS iterations reported & 2040 & 2040 \\
first energy (\kcal{}) & 1347.94 & 1347.94 \\
last energy (\kcal{}) & $-$1102.39 & $-$1102.39 \\
largest difference along the series (\kcal{}) & \multicolumn{2}{c}{0.25} \\
final-structure RMSD (\AA{}) & \multicolumn{2}{c}{0.0045} \\
wall time, Node 26 (s) & 11.5 & 20.8 \\
\bottomrule
\end{tabular}
\end{table}

\paragraph{The record}
The run in Figure~\ref{fig:studio}(b) wrote a record naming Studio 1.2.6, the
engine commit \texttt{73f598a} and the SHA-256 of the WebAssembly build, the
six staged files of the OL21/OL3 terminal deck by digest, the preset and every
parameter it resolved, and the start and end times, with the final structure,
the engine deck, the engine log, the energy series and the trajectory beside
it, each one click away. That record is what a reader who repeats the run
compares against, and it is the same document on every machine that loads the
page, because no part of the computation happens anywhere else.
